\subsection{Relation to the supplied interaction-classification
paper}\label{proof:local_interactions:relation-to-the-supplied-interaction-classification-paper}

\href{https://arxiv.org/abs/2604.13026}{Marwaha and Sud, \emph{A
complexity phase transition at the EPR Hamiltonian}, arXiv:2604.13026},
studies worst-case ground-state energy estimation for positive-weight
sums of a fixed symmetric quantum two-local term. Its input includes the
interaction graph and weights. The paper's classification refers to that
entire problem family; it does not assert that every graph carrying the
same term has identical computational behavior. Its conjectural BPP
assertion for EPR* also must remain labeled conjectural.

The present objects are different: a classical Gibbs law is given, and
AA asks about every globally bounded low-degree polynomial on the
sampled spin coordinates. The counterexamples above use commuting
diagonal interactions and easy ground states. Therefore ground-state
complexity, local Bell-state level ordering, and finite-temperature
polynomial influence are not interchangeable invariants. A useful
analogy is to organize structural hypotheses, not to transfer complexity
labels without a reduction.
